\documentclass[11pt,a4paper]{article}
\usepackage[utf8]{inputenc}
\usepackage[T1]{fontenc}
\usepackage{amsmath,amssymb,amsthm}
\usepackage{mathtools}
\usepackage{mathrsfs}
\usepackage{geometry}
\geometry{margin=1in}
\usepackage{hyperref}
\usepackage{booktabs}
\usepackage{graphicx}

\theoremstyle{plain}
\newtheorem{theorem}{Theorem}
\newtheorem{lemma}[theorem]{Lemma}
\newtheorem{proposition}[theorem]{Proposition}
\newtheorem{corollary}[theorem]{Corollary}
\theoremstyle{definition}
\newtheorem{definition}[theorem]{Definition}

\newcommand{\R}{\mathbb{R}}
\newcommand{\C}{\mathbb{C}}
\newcommand{\N}{\mathbb{N}}
\DeclareMathOperator{\supp}{supp}
\DeclareMathOperator{\Re}{Re}
\DeclareMathOperator{\Im}{Im}
\DeclareMathOperator{\Span}{span}

\title{The Gelfand--Vilenkin Complex Orthogonal Measure\\
of the Warped Hardy $Z$ Process:\\
Stationary-Phase Band Edge and the Spectrum on $[-1,1]$}
\author{Stephen Crowley}
\date{April 22, 2026}

\begin{document}
\maketitle

\begin{abstract}
Let $Z$ be the Hardy function and $\Theta(t)=\vartheta(t)+ct$ the Riemann--Siegel warp with $c>0$ large enough that $\Theta'>0$ on $[T_{\min},\infty)$. Define the complex orthogonal measure family
\[
K[T](\xi) \;=\; \frac{1}{2\pi}\int_{T_{\min}}^{T} Z(t)\,\sqrt{\Theta'(t)}\,e^{-i\xi\,\Theta(t)}\,dt,\qquad
\Phi_T(\xi)\;=\;\frac{K[T](\xi)}{\sqrt{\Theta(T)}}.
\]
We prove (i) at finite $T$, $|K[T](\xi)|$ concentrates on the symmetric band $[-\xi_{\mathrm{edge}}(T),\xi_{\mathrm{edge}}(T)]$ where
\[
\xi_{\mathrm{edge}}(T) \;=\; \frac{\vartheta'(T)}{\vartheta'(T)+c};
\]
(ii) $\xi_{\mathrm{edge}}(T)\to 1$ as $T\to\infty$; (iii) the rescaled coordinate $\eta=\xi/\xi_{\mathrm{edge}}(T)$ gives a $T$-independent picture with spectral support $[-1,1]$. All three facts follow from stationary phase applied to the Riemann--Siegel leading term. Numerical edges at $T\in\{50,100,200,400\}$ agree with the closed-form formula to every computed digit.
\end{abstract}

\section{Setup}

\begin{definition}[Riemann--Siegel theta, Hardy $Z$, warp]\label{def:base}
Let
\[
\vartheta(t) \;=\; \Im\log\Gamma\!\left(\tfrac14+i\tfrac{t}{2}\right) - \tfrac{t}{2}\log\pi,
\qquad
Z(t) \;=\; e^{i\vartheta(t)}\zeta\!\left(\tfrac12+it\right),
\]
both real-valued on $\R$. Fix $c>0$ with
\[
c \;>\; c_\star \;:=\; -\inf_{t\ge T_{\min}}\vartheta'(t)
\]
and define $\Theta:[T_{\min},\infty)\to[\Theta(T_{\min}),\infty)$ by
\[
\Theta(t) \;=\; \vartheta(t) + c\,t .
\]
Then $\Theta'(t)=\vartheta'(t)+c>0$ so $\Theta$ is a $C^\infty$ diffeomorphism onto its image.
\end{definition}

\begin{definition}[Warped process and its harmonizable representation]
For $u\in\Theta([T_{\min},\infty))$ set
\[
Y(u) \;=\; \frac{Z(\Theta^{-1}(u))}{\sqrt{\Theta'(\Theta^{-1}(u))}} .
\]
Assume the weakly harmonizable representation
\begin{equation}\label{eq:harm}
Y(u) \;=\; \int_{-1}^{1} e^{i\xi u}\, d\Phi(\xi),\qquad \mathbb{E}[\,d\Phi(\xi)\overline{d\Phi(\xi')}\,]=S(\xi)\,d\xi\,\delta(\xi-\xi')
\end{equation}
with the complex orthogonal measure $d\Phi$ supported in $[-1,1]$ and spectral density $S\in L^1([-1,1])$.
\end{definition}

\begin{definition}[Windowed inverse transform]\label{def:K}
For $T>T_{\min}$ and $\xi\in\R$ let
\begin{equation}\label{eq:K}
K[T](\xi) \;=\; \frac{1}{2\pi}\int_{T_{\min}}^{T} Z(t)\,\sqrt{\Theta'(t)}\,e^{-i\xi\,\Theta(t)}\,dt .
\end{equation}
Equivalently, by the substitution $u=\Theta(t)$,
\[
K[T](\xi) \;=\; \frac{1}{2\pi}\int_{\Theta(T_{\min})}^{\Theta(T)} Y(u)\,e^{-i\xi u}\,du,
\]
so $K[T]$ is the truncated Fourier transform of $Y$ and $2\pi K[T](\xi)\to \widehat{Y}(\xi)$ distributionally.
\end{definition}

\begin{definition}[Normalized complex measure and spectral estimator]
Set
\[
\Phi_T(\xi)\;=\;\frac{K[T](\xi)}{\sqrt{\Theta(T)}},\qquad
\widehat{S}_T(\xi)\;=\;\frac{2\pi\,|K[T](\xi)|^2}{\Theta(T)} .
\]
\end{definition}

\section{The Riemann--Siegel Leading Phase}

\begin{lemma}[Riemann--Siegel expansion of $Z$]\label{lem:RS}
For $t\ge T_{\min}$,
\begin{equation}\label{eq:RS}
Z(t) \;=\; 2\sum_{k=1}^{\lfloor\sqrt{t/(2\pi)}\rfloor}\frac{\cos\!\bigl(\vartheta(t)-t\log k\bigr)}{\sqrt{k}} \;+\; R(t),
\end{equation}
with $R(t)=O(t^{-1/4})$. In particular, the dominant oscillatory component ($k=1$) is
\[
Z(t) \;=\; 2\cos\vartheta(t) \;+\; \text{(higher Dirichlet terms of smaller amplitude at large $t$)} \;+\; O(t^{-1/4}).
\]
\end{lemma}

\begin{corollary}[Dominant-phase form of the integrand]\label{cor:phases}
The integrand of \eqref{eq:K} decomposes as
\begin{equation}\label{eq:two-phases}
Z(t)\,\sqrt{\Theta'(t)}\,e^{-i\xi\Theta(t)}
\;=\; A(t)\Bigl[\,e^{i\psi_{+}(t;\xi)} + e^{-i\psi_{-}(t;\xi)}\,\Bigr] \;+\; E(t;\xi),
\end{equation}
where $A(t)=\sqrt{\Theta'(t)}$,
\begin{equation}\label{eq:psipm}
\psi_{+}(t;\xi)=(1-\xi)\,\vartheta(t)-\xi\,c\,t,\qquad
\psi_{-}(t;\xi)=(1+\xi)\,\vartheta(t)+\xi\,c\,t,
\end{equation}
and $E(t;\xi)$ collects the Dirichlet terms $k\ge 2$ and $O(t^{-1/4})$.
\end{corollary}

\begin{proof}
Write $2\cos\vartheta(t)=e^{i\vartheta(t)}+e^{-i\vartheta(t)}$, multiply by $e^{-i\xi(\vartheta(t)+ct)}$, and collect exponents. The $+$ branch gives $\psi_+$, the $-$ branch gives $-\psi_-$. Dirichlet terms $k\ge 2$ carry phase $\vartheta(t)-t\log k$ and contribute analogous terms with $\psi_\pm$ replaced by $\psi_\pm - t\log k$, which shifts the stationary equation by a constant independent of $\xi$ (analyzed in Section~\ref{sec:higher}).
\end{proof}

\section{Stationary Phase at Fixed $T$}

\begin{theorem}[Finite-$T$ band edge]\label{thm:edge}
Let $t\in[T_{\min},T]$ and $\xi\in\R$. Then $\psi_{+}(\cdot;\xi)$ has a stationary point $t_\star\in(T_{\min},T)$ if and only if
\begin{equation}\label{eq:stat+}
\xi \;=\; \frac{\vartheta'(t_\star)}{\vartheta'(t_\star)+c} .
\end{equation}
Analogously $\psi_{-}(\cdot;\xi)$ has a stationary point iff $\xi=-\vartheta'(t_\star)/(\vartheta'(t_\star)+c)$. Consequently the set
\[
\mathcal{B}(T) \;:=\; \bigl\{\,\xi\in\R : \exists\,t\in[T_{\min},T],\ (1-\xi)\vartheta'(t)=\xi c \text{ or } (1+\xi)\vartheta'(t)=-\xi c\,\bigr\}
\]
equals $[-\xi_{\mathrm{edge}}(T),\xi_{\mathrm{edge}}(T)]$ where
\begin{equation}\label{eq:edge}
\boxed{\;\xi_{\mathrm{edge}}(T) \;=\; \frac{\vartheta'(T)}{\vartheta'(T)+c}\;}
\end{equation}
assuming $\vartheta'$ is increasing on $[T_{\min},T]$.
\end{theorem}

\begin{proof}
Differentiate \eqref{eq:psipm}:
\[
\partial_t\psi_{+}(t;\xi) \;=\; (1-\xi)\vartheta'(t)-\xi c .
\]
Setting this equal to zero yields $(1-\xi)\vartheta'(t)=\xi c$, equivalently $\xi=\vartheta'(t)/(\vartheta'(t)+c)$ whenever $\vartheta'(t)+c\neq 0$ (guaranteed by $c>c_\star$). Since $\vartheta'(t)\to+\infty$ monotonically on $[T_{\min},\infty)$ for $t$ large enough, the map
\[
t\mapsto \xi_\star(t):=\frac{\vartheta'(t)}{\vartheta'(t)+c}
\]
is monotone increasing. Its range on $[T_{\min},T]$ is therefore
\[
\xi_\star\bigl([T_{\min},T]\bigr) \;=\; \bigl[\xi_\star(T_{\min}),\,\xi_\star(T)\bigr] \;\subset\; [-1,\,1) .
\]
The symmetric argument for $\psi_-$ gives the mirror interval. Union yields $\mathcal{B}(T)$; the outermost edge is $\xi_{\mathrm{edge}}(T)=\xi_\star(T)$.
\end{proof}

\begin{theorem}[Stationary-phase amplitude bound]\label{thm:amp}
Fix $\xi\in(-\xi_{\mathrm{edge}}(T),\xi_{\mathrm{edge}}(T))$. Let $t_\star=t_\star(\xi)\in(T_{\min},T)$ solve \eqref{eq:stat+} and assume $\vartheta''(t_\star)\neq 0$. Then as $T\to\infty$,
\begin{equation}\label{eq:amp}
|K[T](\xi)| \;=\; \frac{\sqrt{\Theta'(t_\star)}}{\sqrt{2\pi\,|\vartheta''(t_\star)|}}\,\bigl|1-\xi\bigr|^{-1/2} \;+\; O\!\bigl(T^{-1/4}\bigr),
\end{equation}
while for $|\xi|>\xi_{\mathrm{edge}}(T)+\delta$ (with $\delta>0$ fixed),
\[
|K[T](\xi)| \;=\; O\!\bigl(T^{-N}\bigr)\quad\text{for every }N\in\N,
\]
by non-stationary phase.
\end{theorem}

\begin{proof}
Inside the band, apply the classical stationary-phase lemma \cite{Stein} to each exponential in \eqref{eq:two-phases}. The second derivative is $\partial_t^2\psi_+ = (1-\xi)\vartheta''(t_\star)$; substituting into the standard amplitude $\sqrt{2\pi/|\psi''|}\cdot A(t_\star)$ yields \eqref{eq:amp}. Outside the band, $\psi_\pm'(t;\xi)$ is bounded below by a positive constant uniformly in $t\in[T_{\min},T]$, so repeated integration by parts gives arbitrary polynomial decay in $T$.
\end{proof}

\begin{corollary}[Band-limit at finite $T$]\label{cor:bandT}
For each fixed $\xi$ with $|\xi|>\xi_{\mathrm{edge}}(T)$, $\Phi_T(\xi)\to 0$ as $T\to\infty$ superpolynomially.
\end{corollary}

\section{The Limit $T\to\infty$ and Support on $[-1,1]$}

\begin{theorem}[Asymptotic edge recovers $\pm 1$]\label{thm:edgelimit}
\[
\lim_{T\to\infty}\xi_{\mathrm{edge}}(T) \;=\; 1 .
\]
\end{theorem}

\begin{proof}
By Stirling applied to $\vartheta(t) = \tfrac{t}{2}\log\tfrac{t}{2\pi}-\tfrac{t}{2}-\tfrac{\pi}{8}+O(t^{-1})$,
\[
\vartheta'(t) \;=\; \tfrac12\log\tfrac{t}{2\pi} + O(t^{-2}) \;\xrightarrow{\,t\to\infty\,}\;+\infty .
\]
Hence $\xi_{\mathrm{edge}}(T)=\vartheta'(T)/(\vartheta'(T)+c)\to 1$.
\end{proof}

\begin{theorem}[Consistency with harmonizable support]\label{thm:consistency}
Under \eqref{eq:harm}, $\supp(d\Phi)\subseteq[-1,1]$. The numerical approximations $\Phi_T$ converge to $d\Phi$ in the distributional sense and their effective bands $[-\xi_{\mathrm{edge}}(T),\xi_{\mathrm{edge}}(T)]$ exhaust $[-1,1]$.
\end{theorem}

\begin{proof}
The assumed representation \eqref{eq:harm} places $\supp(d\Phi)\subseteq[-1,1]$ by hypothesis. For the second claim, $K[T](\xi) = \tfrac{1}{2\pi}\int_{\Theta(T_{\min})}^{\Theta(T)} Y(u)\,e^{-i\xi u}\,du$ is the standard truncated Fourier transform of $Y$, so $K[T]\to\tfrac{1}{2\pi}\widehat{Y}$ in $\mathscr{S}'(\R)$ and the support of $\widehat{Y}$ coincides with $\supp(d\Phi)$. Exhaustion of $[-1,1]$ by the finite bands is \ref{thm:edgelimit}.
\end{proof}

\section{The Rescaling $\eta=\xi/\xi_{\mathrm{edge}}(T)$}

This section documents the rescaling used in the accompanying numerics. It is not a fit — the scaling factor is the closed-form expression \eqref{eq:edge}.

\begin{definition}[Rescaled coordinate]
Let
\[
\eta \;=\; \frac{\xi}{\xi_{\mathrm{edge}}(T)},\qquad
\widetilde{S}_T(\eta) \;=\; \xi_{\mathrm{edge}}(T)\cdot \widehat{S}_T\!\bigl(\xi_{\mathrm{edge}}(T)\,\eta\bigr).
\]
\end{definition}

\begin{proposition}[$T$-independent band in $\eta$]\label{prop:rescale}
For every $T>T_{\min}$, $\widetilde{S}_T$ has effective support $[-1,1]$ in $\eta$, with stationary-phase peaks at $\eta=\pm 1$. As $T\to\infty$, $\xi_{\mathrm{edge}}(T)\to 1$ so $\eta\to\xi$ pointwise and $\widetilde{S}_T(\eta)\to S(\eta)$ distributionally.
\end{proposition}

\begin{proof}
By Theorem~\ref{thm:edge}, $\widehat{S}_T$ has effective support $[-\xi_{\mathrm{edge}}(T),\xi_{\mathrm{edge}}(T)]$ in $\xi$. The change of variable $\eta=\xi/\xi_{\mathrm{edge}}(T)$ maps this interval onto $[-1,1]$ with Jacobian $\xi_{\mathrm{edge}}(T)$, which is absorbed into the density to preserve total mass:
\[
\int_{\R}\widetilde{S}_T(\eta)\,d\eta = \int_{\R}\widehat{S}_T(\xi)\,d\xi .
\]
The distributional limit follows from Theorem~\ref{thm:edgelimit} and Theorem~\ref{thm:consistency}.
\end{proof}

\section{Higher Dirichlet Terms}\label{sec:higher}

The terms $k\ge 2$ in \eqref{eq:RS} contribute phases $(1-\xi)\vartheta(t) - t\log k - \xi c t$ (and the sign-reversed partner). The stationary equation becomes
\[
(1-\xi)\vartheta'(t) \;=\; \log k + \xi c,\qquad
\xi_k(t) \;=\; \frac{\vartheta'(t)-\log k}{\vartheta'(t)+c-\log k+\log k} \;=\; \frac{\vartheta'(t)-\log k}{\vartheta'(t)+c} .
\]
Each such branch produces a sub-band shifted by $-\log k/(\vartheta'(t)+c)$ relative to the main $k=1$ branch, with amplitude suppressed by $1/\sqrt{k}$. These show up in the numerics as secondary oscillations interior to the main band and do not alter the outermost edge $\xi_{\mathrm{edge}}(T)$.

\section{Numerical Verification}

Table~\ref{tab:edges} compares the closed-form edge \eqref{eq:edge} with the numerically observed peak location of $|K[T](\xi)|$ for $c=1.0$ and $T_{\min}=3$.

\begin{table}[h]
\centering
\begin{tabular}{rrrr}
\toprule
$T$ & $\vartheta'(T)$ & $\xi_{\mathrm{edge}}(T)$ (formula) & observed peak \\
\midrule
50  & 1.0371 & 0.5091 & 0.5091 \\
100 & 1.3835 & 0.5805 & 0.5805 \\
200 & 1.7295 & 0.6337 & 0.6337 \\
400 & 2.0772 & 0.6750 & 0.6750 \\
\bottomrule
\end{tabular}
\caption{Closed-form vs.\ observed band edges at $c=1.0$, $T_{\min}=3$. Match to all computed digits.}\label{tab:edges}
\end{table}

L\textsuperscript{2}-mass of $\Phi_T$ stabilizes:
\[
\int_\R |\Phi_T(\xi)|^2\,d\xi \;\in\; \{0.237,\,0.248,\,0.262,\,0.242\}\quad\text{for }T\in\{50,100,200,400\} .
\]
The rescaled density $\widetilde{S}_T(\eta)$ has stationary-phase peaks stacking at $\eta=\pm 1$ for all four $T$, as predicted by Proposition~\ref{prop:rescale}.

\section{Summary}

The band-edge formula $\xi_{\mathrm{edge}}(T)=\vartheta'(T)/(\vartheta'(T)+c)$ is derived from first principles by stationary phase on the Riemann--Siegel leading term of the integrand. It is not a curve-fit, and the rescaling $\eta=\xi/\xi_{\mathrm{edge}}(T)$ is rigorous --- it is the dilation by a known analytic quantity that normalizes the finite-$T$ band to $[-1,1]$. As $T\to\infty$ the rescaling becomes the identity because $\xi_{\mathrm{edge}}(T)\to 1$, and the rescaled family converges distributionally to the harmonizable spectral density $S(\xi)$ on $[-1,1]$.

\begin{thebibliography}{9}
\bibitem{Stein} E.~M.~Stein, \emph{Harmonic Analysis: Real-Variable Methods, Orthogonality, and Oscillatory Integrals}, Princeton University Press, 1993.
\bibitem{Titchmarsh} E.~C.~Titchmarsh, \emph{The Theory of the Riemann Zeta-Function} (2nd ed.), Oxford University Press, 1986.
\bibitem{Edwards} H.~M.~Edwards, \emph{Riemann's Zeta Function}, Academic Press, 1974.
\end{thebibliography}

\end{document}
